
\documentclass[a4paper,12pt]{article}

\usepackage[a4paper,margin=8mm,bmargin=16mm]{geometry}
\usepackage[german,french]{babel}
\usepackage{fontspec}
\usepackage{microtype}
\usepackage{csquotes}

\setmainfont{Alegreya}

\usepackage{paracol}
\newcommand\chunks[2]{%
\begin{leftcolumn*}\begin{otherlanguage}{german}%
{#1}%
\end{otherlanguage}\end{leftcolumn*}%
\begin{rightcolumn}\begin{otherlanguage}{french}%
{#2}%
\end{otherlanguage}\end{rightcolumn}%
}
% https://tex.stackexchange.com/a/505562/295527

\usepackage{luacolor}
\usepackage[soul]{lua-ul}
% https://tex.stackexchange.com/a/747323
%\renewcommand{\hl}{\relax}
%\renewcommand{\hl}[1]{#1}

\usepackage{psvectorian}
\newcommand{\ornamentbreak}{%
  \begin{center}
    \psvectorian[width=2em]{2}\quad\psvectorian[width=2em,mirror]{2}%
  \end{center}%
}

\begin{document}
\thispagestyle{empty}
\setlength{\columnsep}{8mm}
\setlength{\columnseprule}{0.4pt}

\begin{paracol}{2}
  \chunks{\section*{Über die Weiber}}
         {\section*{Sur les femmes}}

  \chunks{\textelp{}}
         {\textelp{}}
  \chunks{Mit den Mädchen hat es die Natur auf Das, was man im dramaturgischen Sinne, einen Knalleffekt nennt, abgesehen, indem sie dieselben, auf wenige Jahre, mit überreichlicher Schönheit, Reiz und Fülle ausstattet, auf Kosten ihrer ganzen übrigen Lebenszeit, damit sie nämlich, während jener Jahre, der Phantasie eines Mannes sich in dem Maße bemächtigen könnten, daß er hingerissen wird, die Sorge für sie auf Zeit Lebens, in irgend einer Form, ehrlich zu übernehmen; zu welchem Schritte ihn zu vermögen, die bloße vernünftige Ueberlegung keine hinlänglich sichere Bürgschaft zu geben schien. Sonach hat die Natur das Weib, eben wie jedes andere ihrer Geschöpfe, mit den Waffen und Werkzeugen ausgerüstet, deren es zur Sicherung seines Daseins bedarf, und auf die Zeit, da es ihrer bedarf; wobei sie denn auch mit ihrer gewöhnlichen Sparsamkeit verfahren ist. Wie nämlich die weibliche Ameise, nach der Begattung, die fortan überflüssigen, ja, für das Brutverhältniß gefährlichen Flügel verliert; so meistens nach einem oder zwei Kindbetten, das Weib seine Schönheit; wahrscheinlich sogar aus dem selben Grunde.}
         {Chez les jeunes filles, la nature semble avoir voulu faire ce qu’en style dramatique on appelle un coup de théâtre ; elle les pare pour quelques années d’une beauté, d’une grâce, d’une perfection extraordinaires, aux dépens de tout le reste de leur vie, afin que pendant ces rapides années d’éclat elles puissent s’emparer fortement de l’imagination d’un homme et l’entraîner à se charger loyalement d’elles d’une manière quelconque. Pour réussir dans cette entreprise la pure réflexion et la raison ne donnaient pas de garantie suffisante. Aussi la nature a-t-elle armé la femme, comme toute autre créature, des armes et des instruments nécessaires pour assurer son existence et seulement pendant le temps indispensable, car la nature en cela agit avec son économie habituelle : de même que la fourmi femelle, après son union avec le mâle, perd les ailes qui lui deviendraient inutiles et même dangereuses pour la période d’incubation, de même aussi la plupart du temps, après deux ou trois couches, la femme perd sa beauté, sans doute pour la même raison.}
  \chunks{Dem entsprechend halten die jungen Mädchen ihre häuslichen, oder gewerblichen Geschäfte, in ihrem Herzen, für Nebensache, wohl gar für bloßen Spaß: als ihren allein ernstlichen Beruf betrachten sie die Liebe, die Eroberungen und was damit in Verbindung steht, wie Toilette, Tanz u. s. w.}
         {De là vient que les jeunes filles regardent généralement les occupations du ménage ou les devoirs de leur état comme des choses accessoires et de pures bagatelles, tandis qu’elles reconnaissent leur véritable vocation dans l’amour, les conquêtes et tout ce qui en dépend, la toilette, la danse, etc.}
  \chunks{Je edler und vollkommener eine Sache ist, desto später und langsamer gelangt sie zur Reife. Der Mann erlangt die Reife seiner Vernunft und Geisteskräfte kaum vor dem acht und zwanzigsten Jahre; das Weib mit dem achtzehnten. Aber es ist auch eine Vernunft danach: eine gar knapp gemessene. Daher bleiben die Weiber ihr Leben lang Kinder, sehn immer nur das Nächste, kleben an der Gegenwart, nehmen den Schein der Dinge für die Sache und ziehn Kleinigkeiten den wichtigen Angelegenheiten vor. Die Vernunft nämlich ist es, vermöge deren der Mensch nicht, wie das Thier, bloß in der Gegenwart lebt, sondern Vergangenheit und Zukunft übersieht und bedenkt; woraus dann seine Vorsicht, seine Sorge und häufige Beklommenheit entspringt. Der Vortheile, wie der Nachtheile, die Dies bringt, ist das Weib, in Folge seiner schwächern Vernunft, weniger theilhaft; vielmehr ist derselbe ein geistiger Myops, indem sein intuitiver Verstand in der Nähe scharf sieht, hingegen einen engen Gesichtskreis hat, in welchen das Entfernte nicht fällt; daher eben alles Abwesende, Vergangene, Künftige, viel schwächer auf die Weiber wirkt, als auf uns, woraus denn auch der bei ihnen viel häufigere und bisweilen an Verrücktheit grenzende Hang zur Verschwendung entspringt. Die Weiber denken in ihrem Herzen, die Bestimmung der Männer sei, Geld zu verdienen, die ihre hingegen, es durchzubringen; wo möglich schon bei Lebzeiten des Mannes, wenigstens aber nach seinem Tode. Schon daß der Mann das Erworbene ihnen zur Haushaltung übergiebt, bestärkt sie in dem Glauben. --- So viele Nachtheile Dies alles zwar mit sich führt, so hat es doch das Gute, daß das Weib mehr in der Gegenwart aufgeht, als wir, und daher diese, wenn sie nur erträglich ist, besser genießen, woraus die dem Weibe eigenthümliche Heiterkeit hervorgeht, welche sie zur Erholung, erforderlichen Falles zum Troste des sorgenbelasteten Mannes eignet.}
         {Plus une chose est noble et accomplie, plus elle se développe lentement et tardivement. La raison et l’intelligence de l’homme n’atteignent guère tout leur développement que vers la vingt-huitième année ; chez la femme, au contraire, la maturité de l’esprit arrive à la dix-huitième année. Aussi n’a-t-elle qu’une raison de dix-huit ans bien strictement mesurée. C’est pour cela que les femmes restent toute leur vie de vrais enfants. Elles ne voient que ce qui est sous leurs yeux, s’attachent au présent, prenant l’apparence pour la réalité et préférant les niaiseries aux choses les plus importantes. Ce qui distingue l’homme de l’animal c’est la raison ; confiné dans le présent, il se reporte vers le passé et songe à l’avenir : de là sa prudence, ses soucis, ses appréhensions fréquentes. La raison débile de la femme ne participe ni à ces avantages, ni à ces inconvénients ; elle est affligée d’une myopie intellectuelle qui lui permet, par une sorte d’intuition, de voir d’une façon pénétrante les choses prochaines ; mais son horizon est borné, ce qui est lointain lui échappe. De là vient que tout ce qui n’est pas immédiat, le passé et l’avenir, agissent plus faiblement sur la femme que sur nous : de là aussi ce penchant bien plus fréquent à la prodigalité, qui parfois touche à la démence. Au fond du cœur les femmes s’imaginent que les hommes sont faits pour gagner de l’argent et les femmes pour le dépenser ; si elles en sont empêchées pendant la vie de leur mari, elles se dédommagent après sa mort. Et ce qui contribue à les confirmer dans cette conviction, c’est que leur mari leur donne l’argent et les charge d’entretenir la maison. --- Tant de côtés défectueux sont pourtant compensés par un avantage : la femme plus absorbée dans le moment présent, pour peu qu’il soit supportable en jouit plus que nous ; de là cet enjouement qui lui est propre et la rend capable de distraire et parfois de consoler l’homme accablé de soucis et de peines.}
  \chunks{In schwierigen Angelegenheiten, nach Weise der alten Germanen, auch die Weiber zu Rathe zu ziehn, ist keineswegs verwerflich: denn ihre Auffassungsweise der Dinge ist von der unsrigen ganz verschieden und zwar besonders dadurch, daß sie gern den kürzesten Weg zum Ziele und überhaupt das zunächst Liegende ins Auge faßt, über welches wir, eben weil es vor unserer Nase liegt, meistens weit hinwegsehn; wo es uns dann Noth thut, darauf zurückgeführt zu werden, um die nahe und einfache Ansicht wieder zu gewinnen. Hiezu kommt, daß die Weiber entschieden nüchterner sind, als wir; wodurch sie in den Dingen nicht mehr sehn, als wirklich da ist; während wir, wenn unsere Leidenschaften erregt sind, leicht das Vorhandene vergrößern, oder Imaginäres hinzufügen.}
         {Dans les circonstances difficiles il ne faut pas dédaigner de faire appel, comme autrefois les Germains, aux conseils des femmes ; car elles ont une manière de concevoir les choses toute différente de la nôtre. Elles vont au but par le chemin le plus court, parce que leurs regards s’attachent, en général, à ce qu’elles ont sous la main. Pour nous, au contraire, notre regard dépasse sans s’y arrêter les choses qui nous crèvent les yeux, et cherche bien au delà ; nous avons besoin d’être ramenés à une manière de voir plus simple et plus rapide. Ajoutez à cela que les femmes ont décidément un esprit plus posé, et ne voient dans les choses que ce qu’il y a réellement ; tandis que, sous le coup de nos passions excitées, nous grossissons les objets, et nous nous peignons des chimères.}
  \chunks{Aus der selben Quelle ist es abzuleiten, daß die Weiber mehr Mitleid und daher mehr Menschenliebe und Theilnahme an Unglücklichen zeigen, als die Männer: hingegen aber im Punkte der Gerechtigkeit, Redlichkeit und Gewissenhaftigkeit, diesen nachstehn. Denn in Folge ihrer schwachen Vernunft übt das Gegenwärtige, Anschauliche, unmittelbar Reale eine Gewalt über sie aus, gegen welche die abstrakten Gedanken, die stehenden Maximen, die festgefaßten Entschlüsse, überhaupt die Rücksicht auf Vergangenheit und Zukunft, auf Abwesendes und Entferntes, selten viel vermögen.}
         {Les mêmes aptitudes natives expliquent la pitié, l’humanité, la sympathie que les femmes témoignent aux malheureux, tandis qu’elles sont inférieures aux hommes en tout ce qui touche à l’équité, à la droiture et à la scrupuleuse probité. À cause de la faiblesse de leur raison, tout ce qui est présent, visible et immédiat, exerce sur elles un empire contre lequel ne sauraient prévaloir ni les abstractions, ni les maximes établies, ni les résolutions énergiques, ni aucune considération du passé ou de l’avenir, de ce qui est éloigné ou absent.}
  \chunks{Demgemäß wird man als den Grundfehler des weiblichen Charakters Ungerechtigkeit finden. Er entsteht zunächst aus dem dargelegten Mangel an Vernünftigkeit und Ueberlegung, wird zudem aber noch dadurch unterstützt, daß sie, als die schwächeren, von der Natur nicht auf die Kraft, sondern auf die List angewiesen sind: daher ihre instinktartige Verschlagenheit und ihr unvertilgbarer Hang zum Lügen. Denn, wie den Löwen mit Klauen und Gebiß, den Elephanten mit Stoßzähnen, den Eber mit Hauern, den Stier mit Hörnern und die Sepia mit der wassertrübenden Tinte, so hat die Natur das Weib mit Verstellungskraft ausgerüstet, zu seinem Schutz und Wehr, und hat alle die Kraft, die sie dem Manne als körperliche Stärke und Vernunft verlieh, dem Weibe in Gestalt jener Gabe zugewendet. Die Verstellung ist ihm demnach angeboren, deshalb auch fast so sehr dem dummen, wie dem klugen Weibe eigen. Von derselben bei jeder Gelegenheit Gebrauch zu machen ist ihm daher so natürlich, wie jenen Thieren, bei Angriff, sogleich ihre Waffen anzuwenden, und empfindet es sich dabei gewissermaßen als seine Rechte gebrauchend. Darum ist ein ganz wahrhaftes, unverstelltes Weib vielleicht unmöglich. Eben deshalb durchschauen sie fremde Verstellung so leicht, daß es nicht rathsam ist, ihnen gegenüber, es damit zu versuchen. --- Aus dem aufgestellten Grundfehler und seinen Beigaben entspringt aber Falschheit, Treulosigkeit, Verrath, Undank u. s. w. Der gerichtlichen Meineide machen Weiber sich viel öfter schuldig, als Männer. Es ließe sich überhaupt in Frage stellen, ob sie zum Eide zuzulassen sind. --- Von Zeit zu Zeit wiederholt sich überall der Fall, daß Damen, denen nichts abgeht, in Kaufmannsläden etwas heimlich einstecken und entwenden.}
         {Aussi l’injustice est-elle le défaut capital des natures féminines. Cela vient du peu de bon sens et de réflexion que nous avons signalé, et ce qui aggrave encore ce défaut, c’est que la nature, en leur refusant la force, leur a donné, pour protéger leur faiblesse, la ruse en partage ; de là leur fourberie instinctive et leur invincible penchant au mensonge. Le lion a ses dents et ses griffes ; l’éléphant, le sanglier ont leurs défenses, le taureau a ses cornes, la seiche a son encre, qui lui sert à brouiller l’eau autour d’elle ; la nature n’a donné à la femme pour se défendre et se protéger que la dissimulation ; cette faculté supplée à la force que l’homme puise dans la vigueur de ses membres et dans sa raison. La dissimulation est innée chez la femme, chez la plus fine, comme chez la plus sotte. Il lui est aussi naturel d’en user en toute occasion qu’à un animal attaqué de se défendre aussitôt avec ses armes naturelles ; et en agissant ainsi, elle a jusqu’à un certain point conscience de ses droits : ce qui fait qu’il est presque impossible de rencontrer une femme absolument véridique et sincère. Et c’est justement pour cela qu’elle pénètre si aisément la dissimulation d’autrui et qu’il n’est pas prudent d’en faire usage avec elle. --- De ce défaut fondamental et de ses conséquences naissent la fausseté, l’infidélité, la trahison, l’ingratitude, etc. Les femmes aussi se parjurent en justice bien plus fréquemment que les hommes, et ce serait une question de savoir si on doit les admettre à prêter serment. --- Il arrive de temps en temps que des dames, à qui rien ne manque, sont surprises dans les magasins en flagrant délit de vol.}
  \chunks{\textelp{}}
         {\textelp{}}
  \chunks{Weil im Grunde die Weiber ganz allein zur Propagation des Geschlechts da sind und ihre Bestimmung hierin aufgeht; so leben sie durchweg mehr in der Gattung, als in den Individuen, nehmen es in ihrem Herzen ernstlicher mit den Angelegenheiten der Gattung, als mit den individuellen. Dies giebt ihrem ganzen Wesen und Treiben einen gewissen Leichtsinn und überhaupt eine von der des Mannes von Grund aus verschiedene Richtung, aus welcher die so häufige und fast normale Uneinigkeit in der Ehe erwächst.}
         {Comme les femmes sont uniquement créées pour la propagation de l’espèce et que toute leur vocation se concentre en ce point, elles vivent plus pour l’espèce que pour les individus, et prennent plus à cœur les intérêts de l’espèce que les intérêts des individus. C’est ce qui donne à tout leur être et à leur conduite une certaine légèreté et des vues opposées à celles de l’homme : telle est l’origine de cette désunion si fréquente dans le mariage, qu’elle en est devenue presque normale.}
  \chunks{Das niedrig gewachsene, schmalschultrige, breithüftige und kurzbeinige Geschlecht das schöne nennen konnte nur der vom Geschlechtstrieb umnebelte männliche Intellekt: in diesem Triebe nämlich steckt seine ganze Schönheit. Mit mehr Fug, als das schöne, könnte man das weibliche Geschlecht das unästhetische nennen. Weder für Musik, noch Poesie, noch bildende Künste haben sie wirklich und wahrhaftig Sinn und Empfänglichkeit; sondern bloße Aefferei, zum Behuf ihrer Gefallsucht, ist es, wenn sie solche affektiren und vorgeben. Das macht, sie sind keines rein objektiven Antheils an irgend etwas fähig, und der Grund hievon ist, denke ich, folgender. Der Mann strebt in Allem eine direkte Herrschaft über die Dinge an, entweder durch Verstehen oder durch Bezwingen derselben. Aber das Weib ist immer und überall auf eine bloße indirekte Herrschaft verwiesen, nämlich mittels des Mannes, als welchen allein es direkt zu beherrschen hat. Darum liegt es in der Weiber Natur, Alles nur als Mittel, den Mann zu gewinnen, anzusehn, und ihr Antheil an irgend etwas Anderem ist immer nur ein simulirter, ein bloßer Umweg, d. h. läuft auf Koketterie und Aefferei hinaus.}
         {Il a fallu que l’intelligence de l’homme fût obscurcie par l’amour pour qu’il ait appelé beau ce sexe de petite taille, aux épaules étroites, aux larges hanches et aux jambes courtes ; toute sa beauté en effet réside dans l’instinct de l’amour. Au lieu de le nommer beau, il eût été plus juste de l’appeler l’inesthétique. Les femmes n’ont ni le sentiment, ni l’intelligence de la musique, pas plus que de la poésie ou des arts plastiques ; ce n’est chez elles que pure singerie, pur prétexte, pure affectation exploitée par leur désir de plaire. Elles sont incapables de prendre une part désintéressée à quoi que ce soit, en voici la raison. L’homme s’efforce en toute chose de dominer directement soit par l’intelligence, soit par la force ; la femme, au contraire, est toujours et partout réduite à une domination absolument indirecte, c’est-à-dire qu’elle n’a de pouvoir que par l’homme, et c’est sur lui seul qu’elle exerce une influence immédiate. En conséquence, la nature porte les femmes à chercher en toutes choses un moyen de conquérir l’homme, et l’intérêt qu’elles semblent prendre aux choses extérieures est toujours une feinte, un détour, c’est-à-dire pure coquetterie et pure singerie.}
  \chunks{Sie sind \textit{sexus sequior} [das geringere Geschlecht], das in jedem Betracht zurückstehende, zweite Geschlecht, dessen Schwäche man demnach schonen soll, aber welchem Ehrfurcht zu bezeugen über die Maßen lächerlich ist und uns in ihren eigenen Augen herabsetzt. Als die Natur das Menschengeschlecht in zwei Hälften spaltete, hat sie den Schnitt nicht gerade durch die Mitte geführt. Bei aller Polarität ist der Unterschied des positiven vom negativen Pol kein bloß qualitativer, sondern zugleich ein quantitativer. --- So haben eben auch die Alten und die orientalischen Völker die Weiber angesehen und danach die ihnen angemessene Stellung viel richtiger erkannt als wir mit unsrer altfranzösischen Galanterie und abgeschmackter Weiberveneration, dieser höchsten Blüthe christlich-germanischer Dummheit, welche nur gedient hat, sie so arrogant und rücksichtslos zu machen, daß man bisweilen an die heiligen Affen in Benares erinnert wird, welche, im Bewußtsein ihrer Heiligkeit und Unverletzlichkeit, sich Alles und Jedes erlaubt halten.}
         {Les femmes sont le \textit{sexus sequior}, le sexe second à tous égards, fait pour se tenir à l’écart et au second plan. Certes, il faut épargner leur faiblesse, mais il est ridicule de leur rendre hommage, et cela même nous dégrade à leurs yeux. La nature, en séparant l’espèce humaine en deux catégories, n’a pas fait les parts égales\textelp{} --- C’est bien ce qu’ont pensé de tout temps les anciens et les peuples de l’Orient ; ils se rendaient mieux compte du rôle qui convient aux femmes, que nous ne le faisons avec notre galanterie à l’ancienne mode française et notre stupide vénération, qui est bien l’épanouissement le plus complet de la sottise germano-chrétienne. Cela n’a servi qu’à les rendre si arrogantes, si impertinentes : parfois elles me font penser aux singes sacrés de Bénarès, qui ont si bien conscience de leur dignité sacro-sainte et de leur inviolabilité, qu’ils se croient tout permis.}

  \chunks{\begin{flushright}
          \large\textsc{Schopenhauer}
          \end{flushright}}
         {}
\end{paracol}
\end{document}
